\documentclass[10pt,a4paper]{amsart}
\usepackage{amsmath,amssymb,amsthm,graphicx}
\usepackage{verbatim}
\newenvironment{reference}{\comment}{\endcomment}
%\newenvironment{reference}{}{}
\newenvironment{slogan}{\comment}{\endcomment}
\newenvironment{history}{\comment}{\endcomment}

% For commutative diagrams we use Xy-pic
\usepackage[all]{xy}

% We use 2cell for 2-commutative diagrams.
\xyoption{2cell}
% Theorem environments.
%
\theoremstyle{plain}
\newtheorem{theorem}[subsection]{Theorem}
\newtheorem{proposition}[subsection]{Proposition}
\newtheorem{lemma}[subsection]{Lemma}

\theoremstyle{definition}
\newtheorem{definition}[subsection]{Definition}
\newtheorem{example}[subsection]{Example}
\newtheorem{exercise}[subsection]{Exercise}
\newtheorem{situation}[subsection]{Situation}

\theoremstyle{remark}
\newtheorem{remark}[subsection]{Remark}
\newtheorem{remarks}[subsection]{Remarks}

\numberwithin{equation}{subsection}

% Macros
%
\def\lim{\mathop{\mathrm{lim}}\nolimits}
\def\colim{\mathop{\mathrm{colim}}\nolimits}
\def\Spec{\mathop{\mathrm{Spec}}}
\def\Hom{\mathop{\mathrm{Hom}}\nolimits}
\def\Ext{\mathop{\mathrm{Ext}}\nolimits}
\def\SheafHom{\mathop{\mathcal{H}\!\mathit{om}}\nolimits}
\def\SheafExt{\mathop{\mathcal{E}\!\mathit{xt}}\nolimits}
\def\Sch{\mathit{Sch}}
\def\Mor{\mathop{\mathrm{Mor}}\nolimits}
\def\Ob{\mathop{\mathrm{Ob}}\nolimits}
\def\Sh{\mathop{\mathit{Sh}}\nolimits}
\def\NL{\mathop{N\!L}\nolimits}
\def\CH{\mathop{\mathrm{CH}}\nolimits}
\def\proetale{{pro\text{-}\acute{e}tale}}
\def\etale{{\acute{e}tale}}
\def\QCoh{\mathit{QCoh}}
\def\Ker{\mathop{\mathrm{Ker}}}
\def\Im{\mathop{\mathrm{Im}}}
\def\Coker{\mathop{\mathrm{Coker}}}
\def\Coim{\mathop{\mathrm{Coim}}}

% Boxtimes
%
\DeclareMathSymbol{\boxtimes}{\mathbin}{AMSa}{"02}

%
% Macros for moduli stacks/spaces
%
\def\QCohstack{\mathcal{QC}\!\mathit{oh}}
\def\Cohstack{\mathcal{C}\!\mathit{oh}}
\def\Spacesstack{\mathcal{S}\!\mathit{paces}}
\def\Quotfunctor{\mathrm{Quot}}
\def\Hilbfunctor{\mathrm{Hilb}}
\def\Curvesstack{\mathcal{C}\!\mathit{urves}}
\def\Polarizedstack{\mathcal{P}\!\mathit{olarized}}
\def\Complexesstack{\mathcal{C}\!\mathit{omplexes}}
% \Pic is the operator that assigns to X its picard group, usage \Pic(X)
% \Picardstack_{X/B} denotes the Picard stack of X over B
% \Picardfunctor_{X/B} denotes the Picard functor of X over B
\def\Pic{\mathop{\mathrm{Pic}}\nolimits}
\def\Picardstack{\mathcal{P}\!\mathit{ic}}
\def\Picardfunctor{\mathrm{Pic}}
\def\Deformationcategory{\mathcal{D}\!\mathit{ef}}
\usepackage{cleveref}
\crefname{originaldefinition}{original definition}{original definitions}
\Crefname{originaldefinition}{Original definition}{Original definitions}
\makeatletter
\newlabel{categories-definition-directed-system}{{0031}{source}{}{sourcepointer-1}{}}
\newlabel{categories-definition-directed-system@cref}{{[originaldefinition][0][]0031}{[1][0][]source}{}{}{}}
\makeatother
\crefname{originalsection}{original section}{original sections}
\Crefname{originalsection}{Original section}{Original sections}
\makeatletter
\newlabel{categories-section-directed-colimits}{{04AX}{source}{}{sourcepointer-2}{}}
\newlabel{categories-section-directed-colimits@cref}{{[originalsection][0][]04AX}{[1][0][]source}{}{}{}}
\makeatother
\crefname{originalreference}{original reference}{original references}
\Crefname{originalreference}{Original reference}{Original references}
\makeatletter
\newlabel{lemma-finite-finitely-generated}{{0BU1}{source}{}{sourcepointer-3}{}}
\newlabel{lemma-finite-finitely-generated@cref}{{[originalreference][0][]0BU1}{[1][0][]source}{}{}{}}
\makeatother
\crefname{originalreference}{original reference}{original references}
\Crefname{originalreference}{Original reference}{Original references}
\makeatletter
\newlabel{lemma-galois-infinite}{{0BMK}{source}{}{sourcepointer-4}{}}
\newlabel{lemma-galois-infinite@cref}{{[originalreference][0][]0BMK}{[1][0][]source}{}{}{}}
\makeatother
\crefname{originalreference}{original reference}{original references}
\Crefname{originalreference}{Original reference}{Original references}
\makeatletter
\newlabel{lemma-normal-closure-inside-normal}{{0BMG}{source}{}{sourcepointer-5}{}}
\newlabel{lemma-normal-closure-inside-normal@cref}{{[originalreference][0][]0BMG}{[1][0][]source}{}{}{}}
\makeatother
\crefname{originalreference}{original reference}{original references}
\Crefname{originalreference}{Original reference}{Original references}
\makeatletter
\newlabel{lemma-ses-galois}{{0BMH}{source}{}{sourcepointer-6}{}}
\newlabel{lemma-ses-galois@cref}{{[originalreference][0][]0BMH}{[1][0][]source}{}{}{}}
\makeatother
\crefname{originallemma}{original lemma}{original lemmas}
\Crefname{originallemma}{Original lemma}{Original lemmas}
\makeatletter
\newlabel{topology-lemma-bijective-map}{{08YE}{source}{}{sourcepointer-7}{}}
\newlabel{topology-lemma-bijective-map@cref}{{[originallemma][0][]08YE}{[1][0][]source}{}{}{}}
\makeatother
\crefname{originallemma}{original lemma}{original lemmas}
\Crefname{originallemma}{Original lemma}{Original lemmas}
\makeatletter
\newlabel{topology-lemma-topological-group-limits}{{0B20}{source}{}{sourcepointer-8}{}}
\newlabel{topology-lemma-topological-group-limits@cref}{{[originallemma][0][]0B20}{[1][0][]source}{}{}{}}
\makeatother
\title{The Galois group of an infinite Galois extension is the inverse limit of the finite Galois subextension groups}
\author{The Stacks Project Authors}
\date{stacks-project:fields.tex:local-snapshot-20261012; GNU FDL 1.2 or later}
\begin{document}
\maketitle
\noindent\textit{Editorial note.} All delivered mathematical text is the author's own, copied from the pinned Stacks Project chapter file fields.tex; no statement or proof step was rewritten, reordered or removed.\par\medskip
\noindent\textit{Editorial note.} The counted proof is the author's complete argument: Each lambda is realized as L\_lambda = K(S\_i) for a finite set S\_i (lemma-finite-finitely-generated), which is what makes the directedness argument by union of finite generating sets work.\par\medskip
\noindent\textit{Editorial note.} The label lemma-infinite-galois-limit names the counted statement, and the proof environment immediately following it in the source is delivered verbatim as its proof.\par\medskip
\noindent\textit{Editorial note.} The counted statement is carried by the author's own lemma environment in the chapter source.\par\medskip
\noindent\textit{Allowed standard prerequisites (editorial).}\begin{itemize}
\item Same chapter: lemma-normal-closure-inside-normal (finite Galois E containing any finite K(S) inside L), lemma-finite-finitely-generated (finite extensions are finitely generated), lemma-ses-galois (restriction maps G\_lambda -> G\_lambda' surjective in towers of finite Galois extensions), lemma-galois-infinite (canonical surjective continuous restriction Gal(L/K) -> Gal(L\_lambda/K))
\item Categories, Definition categories-definition-directed-system and Section categories-section-directed-colimits: directed posets and the description of filtered colimits of fields as unions; Topology, Lemma topology-lemma-topological-group-limits (limits of topological groups commute with the forgetful functor) and topology-lemma-bijective-map (a bijective continuous map of profinite spaces is a homeomorphism); cited as standard facts of sibling chapters
\end{itemize}\medskip
\ifdefined\hypertarget\hypertarget{sourcepointer-1}{}\fi
\noindent\textit{Source reference.} Original definition 0031: sources/stacks-chow/tags.txt, source lines 118--118. Prerequisite pointer invoked at 3234; it names a standard fact of a sibling chapter that the counted argument uses, and every step of the counted proof built on it is retained verbatim in the delivered body.\par\medskip
\ifdefined\hypertarget\hypertarget{sourcepointer-2}{}\fi
\noindent\textit{Source reference.} Original section 04AX: sources/stacks-chow/tags.txt, source lines 5632--5632. Prerequisite pointer invoked at 3238; it names a standard fact of a sibling chapter that the counted argument uses, and every step of the counted proof built on it is retained verbatim in the delivered body.\par\medskip
\ifdefined\hypertarget\hypertarget{sourcepointer-3}{}\fi
\noindent\textit{Source reference.} Original reference 0BU1: sources/stacks-fields/pinned-original-source.tex, source lines 707--707. Prerequisite pointer invoked at 3229; it names a standard fact of the same chapter that the counted argument uses, and every step of the counted proof built on it is retained verbatim in the delivered body.\par\medskip
\ifdefined\hypertarget\hypertarget{sourcepointer-4}{}\fi
\noindent\textit{Source reference.} Original reference 0BMK: sources/stacks-fields/pinned-original-source.tex, source lines 3168--3168. Prerequisite pointer invoked at 3249; it names a standard fact of the same chapter that the counted argument uses, and every step of the counted proof built on it is retained verbatim in the delivered body.\par\medskip
\ifdefined\hypertarget\hypertarget{sourcepointer-5}{}\fi
\noindent\textit{Source reference.} Original reference 0BMG: sources/stacks-fields/pinned-original-source.tex, source lines 2351--2351. Prerequisite pointer invoked at 3224; it names a standard fact of the same chapter that the counted argument uses, and every step of the counted proof built on it is retained verbatim in the delivered body.\par\medskip
\ifdefined\hypertarget\hypertarget{sourcepointer-6}{}\fi
\noindent\textit{Source reference.} Original reference 0BMH: sources/stacks-fields/pinned-original-source.tex, source lines 3042--3042. Prerequisite pointer invoked at 3245; it names a standard fact of the same chapter that the counted argument uses, and every step of the counted proof built on it is retained verbatim in the delivered body.\par\medskip
\ifdefined\hypertarget\hypertarget{sourcepointer-7}{}\fi
\noindent\textit{Source reference.} Original lemma 08YE: sources/stacks-chow/tags.txt, source lines 11337--11337. Prerequisite pointer invoked at 3267; it names a standard fact of a sibling chapter that the counted argument uses, and every step of the counted proof built on it is retained verbatim in the delivered body.\par\medskip
\ifdefined\hypertarget\hypertarget{sourcepointer-8}{}\fi
\noindent\textit{Source reference.} Original lemma 0B20: sources/stacks-chow/tags.txt, source lines 13913--13913. Prerequisite pointer invoked at 3261; it names a standard fact of a sibling chapter that the counted argument uses, and every step of the counted proof built on it is retained verbatim in the delivered body.\par\medskip
\begin{lemma}
\label{lemma-infinite-galois-limit}
Let $L/K$ be a Galois extension with Galois group $G$.
Let $\Lambda$ be the set of finite Galois subextensions,
i.e., $\lambda \in \Lambda$ corresponds to $L/L_\lambda/K$
with $L_\lambda/K$ finite Galois with Galois group $G_\lambda$.
Define a partial ordering on $\Lambda$ by the rule
$\lambda \geq \lambda'$ if and only if
$L_\lambda \supset L_{\lambda'}$. Then
\begin{enumerate}
\item $\Lambda$ is a directed partially ordered set,
\item $L_\lambda$ is a system of $K$-extensions over $\Lambda$
and $L = \colim L_\lambda$,
\item $G_\lambda$ is an inverse system of finite groups over $\Lambda$,
the transition maps are surjective, and
$$
G = \lim_{\lambda \in \Lambda} G_\lambda
$$
as a profinite group, and
\item each of the projections $G \to G_\lambda$ is continuous and surjective.
\end{enumerate}
\end{lemma}

\begin{proof}
Every subfield of $L$ containing $K$ is separable over $K$
(follows immediately from the definition). Let $S \subset L$
be a finite subset. Then $K(S)/K$ is finite and there exists
a tower $L/E/K(S)/K$ such that $E/K$ is finite Galois, see
Lemma \ref{lemma-normal-closure-inside-normal}.
Hence $E = L_\lambda$ for some $\lambda \in \Lambda$.
This certainly implies the set $\Lambda$ is not empty.
Also, given $\lambda_1, \lambda_2 \in \Lambda$ we can
write $L_{\lambda_i} = K(S_i)$ for finite sets
$S_1, S_2 \subset L$ (Lemma \ref{lemma-finite-finitely-generated}).
Then there exists a $\lambda \in \Lambda$ such that
$K(S_1 \cup S_2) \subset L_\lambda$. Hence
$\lambda \geq \lambda_1, \lambda_2$ and
$\Lambda$ is directed (Categories, Definition
\ref{categories-definition-directed-system}).
Finally, since every element in $L$ is contained in
$L_\lambda$ for some $\lambda \in \Lambda$, it follows
from the description of filtered colimits in
Categories, Section \ref{categories-section-directed-colimits}
that $\colim L_\lambda = L$.

\medskip\noindent
If $\lambda \geq \lambda'$ in $\Lambda$, then we obtain a
canonical surjective map $G_\lambda \to G_{\lambda'}$,
$\sigma \mapsto \sigma|_{L_{\lambda'}}$
by Lemma \ref{lemma-ses-galois}. Thus we get an inverse
system of finite groups with surjective transition maps.

\medskip\noindent
Recall that $G = \text{Aut}(L/K)$. By Lemma \ref{lemma-galois-infinite}
the restriction $\sigma|_{L_\lambda}$ of a $\sigma \in G$ to $L_\lambda$
is an element of $G_\lambda$. Moreover, this procedure gives a continuous
surjection $G \to G_\lambda$. Since the transition mappings
in the inverse system of $G_\lambda$ are given by restriction
also, it is clear that we obtain a canonical continuous map
$$
G \longrightarrow \lim_{\lambda \in \Lambda} G_\lambda
$$
Continuity by definition of limits in the category of topological
groups; recall that these limits commute with the forgetful functor
to the categories of sets and topological spaces by
Topology, Lemma \ref{topology-lemma-topological-group-limits}.
On the other hand, since $L = \colim L_\lambda$ it is clear
that any element of the inverse limit (viewed as a set) defines an
automorphism of $L$. Thus the map is bijective. Since the topology
on both sides is profinite, and since a bijective continuous map
of profinite spaces is a homeomorphism
(Topology, Lemma \ref{topology-lemma-bijective-map}), the proof is complete.
\end{proof}

\end{document}
